\section{Background}
\label{sec:background}

\subsection{State Estimation: Filtering and Smoothing}
\label{subsec:state-estimation}

For planar motion, let
$\mathbf{x}_k=(p_x,p_y,v_x,v_y)^{\mathsf T}$. A linear state-space model
describes its evolution and observation:
\begin{equation}
  \mathbf{x}_k=\mathbf{F}_k\mathbf{x}_{k-1}+\mathbf{w}_k,
  \qquad
  \mathbf{z}_k=\mathbf{H}_k\mathbf{x}_k+\mathbf{v}_k,
  \label{eq:state-space}
\end{equation}
Here $\mathbf{w}_k$ and $\mathbf{v}_k$ have covariances $\mathbf{Q}_k$
and $\mathbf{R}_k$. Under constant velocity, process noise covers
accelerations~\cite{barshalom2001estimation}. The Kalman filter uses
measurements through $k$~\cite{kalman1960filtering}; Rauch--Tung--Striebel
(RTS) smoothing also uses later measurements~\cite{rauch1965smoothing}.
An impact calls for a new motion segment because velocity can change
abruptly.

\subsection{Multi-Object Tracking and Data Association}
\label{subsec:multi-object-tracking}

Tracking associates detections with trajectories and manages their
initiation and termination~\cite{barshalom2001estimation}. For a
predicted state $\hat{\mathbf{x}}$ with covariance $\mathbf{P}$, the
innovation $\mathbf{r}=\mathbf{z}-\mathbf{H}\hat{\mathbf{x}}$ has covariance
$\mathbf{S}=\mathbf{H}\mathbf{P}\mathbf{H}^{\mathsf T}+\mathbf{R}$.
Its squared Mahalanobis distance is
\begin{equation}
  d^2=\mathbf{r}^{\mathsf T}\mathbf{S}^{-1}\mathbf{r}.
  \label{eq:association-distance}
\end{equation}
A gate rejects incompatible candidates; passing it does not determine
a unique association.

Radar returns arise from visible parts of an extended vehicle
body~\cite{hammarstrand2012extended,granstrom2017extended}. Doppler
measures line-of-sight velocity, while a tracked return need not
coincide with the vehicle centre. Different observers can therefore
track different surfaces of the same car.

With several observers, track-to-track association must determine whether
histories from different sources describe the same physical object, a
problem distinct from fusing their estimates~\cite{chong2000association}.

\subsection{Traffic Conflict Measures}
\label{subsec:traffic-conflicts}

Time to collision (TTC) is the time until predicted contact under stated
motion assumptions~\cite{hayward1972near-miss}. For straight-path
following at constant speeds with $v_f>v_l$, TTC and time gap are
\begin{equation}
  \mathrm{TTC}=\frac{d}{v_f-v_l},
  \qquad
  \tau_{\mathrm{gap}}=\frac{d}{v_f},
  \label{eq:ttc-time-gap}
\end{equation}
where $d>0$ is bumper-to-bumper clearance. At equal speeds, a short
gap can coexist with infinite constant-speed TTC~\cite{vogel2003headway}.

For cut-ins and crossings, prediction must account for both paths and
footprints. Braking-based measures also assume a response delay and
available deceleration~\cite{westhofen2023criticality}.

The interpretation of response timing likewise depends on context.
Naturalistic studies show that braking onset varies with the evolving
kinematic urgency of an encounter, rather than following a universal
delay after a single hazard onset~\cite{markkula2016braking}. An interval
from a conflict observable in sensor data to brake application therefore
describes an observed braking delay; it does not establish when the driver
perceived the hazard or whether the driver was distracted.
